\section{Introduction}

Therapeutic vaccination against tumor neoantigens has emerged as one of the most promising strategies in precision oncology, offering the possibility of eliciting durable, tumor-specific cytotoxic T lymphocyte responses while sparing healthy tissue \citep{Castillo2024}. Unlike shared tumor-associated antigens, neoantigens arise from somatic mutations unique to an individual tumor and are not subject to central immune tolerance, making them attractive targets for personalized immunotherapy \citep{Reyes2023}. Messenger RNA (mRNA) has become the delivery modality of choice for encoding these antigens, combining rapid, sequence-agnostic manufacturing with the ability to encode multiple epitopes in a single construct \citep{Anand2025}.

The clinical success of mRNA vaccines nevertheless hinges on efficient delivery to antigen-presenting cells, and in particular on cross-presentation: the process by which dendritic cells (DCs) route exogenously acquired antigen into the MHC class~I pathway to prime naive CD8$^{+}$ T cells \citep{Lindgren2022}. Conventional ionizable lipid nanoparticles (LNPs) achieve efficient endosomal uptake by DCs, but most internalized cargo remains trapped within maturing endosomes and is degraded before reaching the cytosol, where proteasomal processing and MHC-I loading must occur \citep{Vertex2021}. Fewer than $2\%$ of endocytosed LNP cargo is estimated to achieve productive cytosolic release under standard formulations, a major bottleneck for cross-presentation efficiency \citep{Nimbus2023}. Strategies explored to improve escape, including pH-responsive lipids and fusogenic peptide conjugates, often introduce cytotoxicity or rely on triggers difficult to translate to systemic administration \citep{Meridian2024}.

We reasoned that a surface-functionalization strategy targeting DC-specific uptake and trafficking, rather than a bulk lipid-chemistry change, could improve cytosolic delivery while preserving formulation simplicity, building on evidence that C-type lectin receptor engagement can route nanoparticle cargo through DC-selective uptake pathways \citep{Falk2020}. Here we describe Synapse-LNP, an ionizable lipid nanoparticle platform surface-functionalized with a DC-targeting glycopeptide ligand (the ``Synapse'' moiety) designed to engage C-type lectin receptors on conventional type~1 DCs and promote early endosomal escape. Our objectives were to:
\begin{enumerate}[label=\sffamily\arabic*.]
\item Formulate and physicochemically characterize Synapse-LNPs encapsulating mRNA encoding validated B16F10 melanoma neoantigens.
\item Quantify DC uptake, maturation, and cross-presentation efficiency in vitro relative to unmodified LNPs and free peptide.
\item Evaluate neoantigen-specific CD8$^{+}$ T cell responses, tumor growth control, and survival in an immunocompetent murine melanoma model.
\end{enumerate}
Together, these experiments test whether targeted endosomal-escape engineering translates into a measurable gain in vaccine-induced antitumor immunity.
